\documentclass[11pt,a4paper,openany]{book}
\usepackage[utf8]{inputenc}
\usepackage[T1]{fontenc}
\usepackage[french]{babel}
\usepackage{amsmath,amssymb,amsthm}
\usepackage{geometry}
\geometry{margin=2.5cm}

\newtheorem{theoreme}{Théorème}[chapter]
\newtheorem{definition}{Définition}[chapter]

\title{\Huge \textbf{LA MACHINE v22.0}\\[0.4cm]
\Large Arithmétique des Roues Primorielles, Cryptographie Post-Quantique et Preuves Formelles en Lean 4}
\author{\textbf{Mohamed Ghali} — ORCID : 0009-0006-3364-5341}
\date{Novembre 2026}

\begin{document}
\frontmatter
\maketitle
\tableofcontents
\mainmatter


\chapter*{Préface}
Cet ouvrage rassemble les fondements théoriques, algorithmiques et formels de \textbf{La Machine v22.0}. L'intégralité des résultats mathématiques présentés est accompagnée de sa certification mécanique en \textbf{Lean 4 pur} et d'une implémentation reproductible en \textbf{Python}.

\chapter{Roues Primorielles et Mode Rafale}
\section{Filtrage géométrique préalable}
Pour tout nombre premier $p_r$, le socle primoriel $S_r = \prod_{p \le p_r} p$ projette la recherche sur le groupe des unités $(\mathbb{Z}/S_r\mathbb{Z})^\times$, réduisant l'espace des candidats à la fraction d'Euler $\varphi(S_r)/S_r$.

\begin{theoreme}[Symétrie Miroir de la Roue Primorielle]
Pour tout socle $S$ et tout résidu $x < S$, si $\gcd(x, S) = 1$, alors son symétrique $S - x$ vérifie également $\gcd(S - x, S) = 1$.
\end{theoreme}

\section{Le Mode Rafale (\texttt{next\_prime\_burst\_mode})}
Grâce à la troncature d'échelle (\texttt{scaleTrunc}), lorsque la fenêtre $W < S$, l'exploration se limite à un seul tour de roue et élimine plus de $88\,\%$ des entiers composés avant le test de Miller-Rabin déterministe.
\chapter{Cryptographie Post-Quantique et Corps Premiers NTT}

\section{Certification des corps premiers de Proth}
Les protocoles post-quantiques standardisés (\textbf{ML-KEM}, \textbf{Falcon-1024}) ainsi que les preuves ZK-STARKs (\textbf{BabyBear}, \textbf{KoalaBear}) reposent sur des corps premiers de la forme $p = k \cdot 2^n + 1$ à faible poids de Hamming.

\begin{theoreme}[Racine primitive NTT et demi-exposant]
Pour tout premier de Proth $p = k \cdot 2^n + 1$ ($n \ge 1$) et tout témoin $a$ vérifiant $a^{(p-1)/2} \equiv -1 \pmod{p}$, l'élément $\omega_{2^n} = a^k \bmod p$ satisfait $\omega_{2^n}^{2^{n-1}} \equiv -1 \pmod{p}$ et engendre un sous-groupe multiplicatif d'ordre exact $2^n$.
\end{theoreme}

\section{Réduction paresseuse K-RED et extension $\mathbb{F}_{p^2}$}
L'identité algébrique $k \cdot (H \cdot 2^n + L) + H = H \cdot p + k \cdot L$ élimine toute division matérielle lors des papillons NTT (\texttt{prothKRedTwoLayers}). L'extension quadratique $\mathbb{F}_{p^2} = \mathbb{F}_p[u]/(u^2 - \beta)$ par multiplication de Karatsuba porte ensuite la sécurité cryptographique à $128$~bits.
\end{document}
